\subsection{Minimal moves and the geometry of finite strands}
\label{sec:examples}
We now specialize the preceding construction to its combinatorial
description. A monomial is a vertex, a chosen binary relation gives
an edge, and a syzygy records a relation among those edges. The
following order fixes the quadratic generators and their orientations.
Higher cells will be specified by their boundaries in individual
strands; no general cellular resolution is assumed.
\subsubsection{The product order}
In each strand we compare monomials of the same bidegree $(n,d)$.
Write them as $x_{\boldsymbol a}=x_{a_1}\cdots x_{a_d}$ and
$x_{\boldsymbol b}=x_{b_1}\cdots x_{b_d}$, so that
$\sum_j a_j=\sum_j b_j=n$. On the odd primes our order is simply
\[
 x_{\boldsymbol a}\prec x_{\boldsymbol b}
 \quad\Longleftrightarrow\quad
 a_1\cdots a_d<b_1\cdots b_d.
\]
Unique factorization distinguishes all the monomials in the fiber.

The connection with \emph{majorization} explains why balanced pairs
give the largest products. List both partitions in decreasing order.
Then $\boldsymbol a$ majorizes $\boldsymbol b$ if
$\sum_{j=1}^s a_j\ge\sum_{j=1}^s b_j$ for each $1\le s<d$.
For distinct partitions of positive labels this implies
$a_1\cdots a_d<b_1\cdots b_d$, by strict concavity of $\log$ and
the majorization inequality \cite[Ch.~3]{MOA11}.
For $d=2$ the larger of the two parts determines majorization.
For $d=3$ the product also compares partitions that majorization
leaves incomparable. In $R_{31,3}$, for instance, the first two
partial sums for $(13,13,5)$ are $13,26$, whereas those for
$(17,7,7)$ are $17,24$. Their products are $845$ and $833$.

For an arbitrary bounded-below support $\A$, put
$c=\max\{0,2-\min\A\}$ and compare first by
$P_c(x^\alpha)=\prod_{a\in\A}(a+c)^{\alpha(a)}$.
The same majorization argument applies to $\log(a+c)$.
If two products agree, take the largest label $a_0$ at which the
exponents differ and put $x^\alpha\prec x^\beta$ if
$\alpha(a_0)<\beta(a_0)$. For example, $x_2x_6^2$ and $x_3^2x_8$
on $\A=\{2,3,6,8\}$ have the same sum, monomial degree and product.
The comparison by $\log P_c$ is a weight order with weights
$\log(a+c)$ \cite[Definition~1.2 and Section~6]{Vietri02}.

\begin{proposition}[Product monomial order]\label{moves:order}
The same comparison on all monomials is a monomial well-order on $R$.
With $c$ fixed, it restricts to each polynomial subring on a subset
of the variables.
\end{proposition}
\begin{proof}
The product comparison followed by the stated lexicographic tie-break
is a total order, with $1$ its least element. Products multiply, and
adding the same exponent vector preserves the largest differing label
and its exponent comparison. Hence multiplying by any monomial
preserves strict inequalities.

For $N\ge1$, only finitely many monomials satisfy $P_c\le N$.
Every factor satisfies $2\le a+c\le N$, and the number of factors
is at most $\log_2 N$. Every nonempty set of monomials therefore has a
least element, obtained by taking its least product score and then
the least of the finitely many ties. Restriction changes none of
these comparisons.
\end{proof}

\subsubsection{Minimal binary Goldbach moves}
For every represented sum $n$, let
\[
 \mathcal M_{n,2}=\{x_ax_{a'}:a,a'\in\A,\ a\le a',\ a+a'=n\}.
\]
Choose $b\le b'$ in $\A$ with $b+b'=n$ and $b'-b$ as small as
possible. We call $u_n=x_bx_{b'}$ the \emph{anchor} of this fiber.
If $n/2\in\A$, the anchor is the square $x_{n/2}^2$.

\begin{proposition}[Anchors and minimal generators]\label{moves:anchors}
The anchor $u_n$ is the unique largest monomial in $\mathcal M_{n,2}$
for the product order, independently of the shift $c$. The binomials
\[
 u_n-x_ax_{a'},\qquad
 x_ax_{a'}\in\mathcal M_{n,2}\setminus\{u_n\},
\]
as $n$ varies, form a minimal bihomogeneous generating set of $I$
over every field. If a quadratic fiber is nonempty in $R^{<m}$, its
anchor already belongs to $R^{<m}$. Consequently these chosen
generating sets are nested under variable truncation.
\end{proposition}
\begin{proof}
For $a+a'=n$ and $a\le a'$,
\[
 (a+c)(a'+c)=\frac{(n+2c)^2-(a'-a)^2}{4}.
\]
The sum and the gap determine the pair. Thus $u_n$ is the unique
product maximum for every $c$. For any other pair we have
$a<b\le b'<a'$: decreasing the gap at fixed sum raises the smaller
label and lowers the larger one. Hence if $a'\le m$, both $b$ and $b'$
are at most $m$. This proves the assertion about cutoffs.

For any two monomials $v,w\in\mathcal M_{n,2}$, their difference
is $(u_n-w)-(u_n-v)$. Within a fiber, the displayed binomials are
linearly independent because each nonanchor monomial appears in just one.
Different fibers have different Goldbach degrees. Since all generators
have monomial degree two and $\m I$ has monomial degree at least
three, their classes form a basis of $I/\m I$. This proves
minimality. The fixed anchors give the asserted nesting.
\end{proof}

\begin{definition}[Minimal binary Goldbach move]\label{moves:minimal}
A \emph{minimal binary Goldbach move} is a replacement
\[
 x^\alpha x_ax_{a'}\longrightarrow x^\alpha u_n,\qquad
 x_ax_{a'}\in\mathcal M_{n,2}\setminus\{u_n\},
\]
where $x^\alpha$ is the fixed monomial factor. The arrow points
from the nonanchor multiple to the anchor multiple, and hence
increases the product order. We also allow traversal in the opposite
direction when discussing undirected paths.
\end{definition}

Here ``minimal'' refers to the chosen minimal generating set of the
ideal. When we wish to specify a quadratic generator by the monomial
$v=x_ax_{a'}$ it replaces, we write $\left.e\right|_v$. Thus
$\partial(\left.e\right|_v)=u_n-v$ for $n=a+a'$.

\begin{proposition}[The vertices and edges of a strand]\label{moves:graph}
Choose these quadratic generators as the basis of $F_1$.
For every bidegree $(n,d)$,
the map $F_{1,n,d}\to F_{0,n,d}$ is the oriented incidence map
of the graph whose vertices are the monomials $x^\alpha$ of Goldbach degree
$n$ and monomial degree $d$
and whose edges are the minimal binary Goldbach moves in that fiber.
Its connected components are those of $\Gamma_{n,d}$.
\end{proposition}
\begin{proof}
Since $F_{0,n,d}=R_{n,d}$, its basis consists of the stated monomials.
The basis of $F_{1,n,d}$ consists of the multiples
$x^\alpha(\left.e\right|_v)$ of bidegree $(n,d)$. For $v=x_ax_{a'}$,
\[
 \partial\bigl(x^\alpha(\left.e\right|_v)\bigr)
     =x^\alpha u_{a+a'}-x^\alpha x_ax_{a'}.
\]
This is exactly the incidence boundary of the indicated edge.
Any elementary binary move between two multiples of pairs of sum
$a+a'$ with the same fixed factor $x^\alpha$ passes through $x^\alpha u_{a+a'}$
using at most two star edges. Thus the chosen edges connect the
vertices in each component of $\Gamma_{n,d}$.
\end{proof}

Along an arrow $x^\alpha x_ax_{a'}\to x^\alpha u_n$, the pair $(a',a)$ majorizes
$(b',b)$. Appending the labels of $x^\alpha$ to both pairs
preserves majorization: the sum of the $j$ largest labels is the
maximum over all ways to choose $j$ labels from the pair and from $x^\alpha$.
Thus the source partition majorizes the target along every directed path.
Every arrow strictly increases the product order, so the move graph
has no directed cycles. Its reachability relation is a partial order,
but the graph need not be its Hasse diagram: a chosen edge may also
be traversed by a longer directed path. Nor do these quadratic
generators have to be a Gr\"obner basis. The product order selects
and orients the edges; choosing higher syzygies and realizing their
boundaries by cells is a further problem.
Majorization alone does not create an edge: $\Gamma_{n,d}$ allows all
binary replacements, while the chosen graph retains only anchor moves.

\begin{proposition}[Maximal representatives]\label{moves:maximal-representatives}
Every connected component $C$ of $\Gamma_{n,d}$ has a unique greatest
monomial $m_C$ in the product order, called its \emph{component maximum}.
The classes of these monomials form a basis of $G_{n,d}$.
From any vertex, successive product-increasing minimal binary moves
terminate at a \emph{one-move local maximum}: a monomial from which
every nontrivial binary move strictly decreases $P_c$.
\end{proposition}
\begin{proof}
The fiber is finite, so each component has a greatest element in the
total product order, and a strictly increasing sequence of moves
terminates. At a terminal vertex every pair of factors is its sum's
anchor; otherwise replacing that pair by its anchor would increase
the product. Every nontrivial replacement of an anchor decreases
$P_c$. The basis assertion is Proposition~\ref{setup:components}.
\end{proof}

\begin{example}[Local maxima need not be unique]\label{moves:local-maxima}
On the odd primes, the minimal moves in bidegree $(21,3)$ include
\[
 x_5^2x_{11}\;\longleftrightarrow\;
 x_3x_7x_{11}\;\longleftrightarrow\;x_7^3,
 \qquad 275\;>\;231\;<\;343.
\]
The displayed numbers are the products of the labels. Both endpoints
are one-move local maxima: their quadratic factors are the anchors
$u_{10}=x_5^2$, $u_{16}=x_5x_{11}$, and $u_{14}=x_7^2$.
The component maximum is $x_7^3$, since the product of three positive
labels with sum $21$ is at most $7^3$. Thus reaching the component
maximum from a local maximum can require a product-decreasing move.
\end{example}



\subsubsection{Illustrated minimal-resolution strands}
For the examples in this subsection, let $\A=\Podd$ be the positive
odd primes and let $G$ be the full binary quotient. We now describe
the $\kk$-complex $F_{n,d}$ for several choices of
$(n,d)$. The vertices are labeled by the monomials $x^\alpha\in R_{n,d}$,
the edges correspond to the chosen quadratic Goldbach relations,
and the two-dimensional cells record relations among these moves.
We use the maximum-product anchor in each quadratic fiber. A
quadratic basis element has boundary equal to the anchor
minus the nonanchor monomial. Arrows in the following pictures indicate this
cellular boundary orientation; they point toward increasing product.

We label a generator of Goldbach degree $n$ and monomial degree $d$
by $e_{n,d}$, using $e_{n,d}',e_{n,d}'',\ldots$ when several occur.
The two subscripts record its bidegree. In the quadratic fiber we
label the generators $\left.e\right|_v$ in order of increasing
$P_c(u_n)-P_c(v)$, beginning with $e_{n,2}$. These gaps are distinct.
Thus each additional prime indicates a smaller second monomial $v$
in the product order. In higher degrees the primes distinguish the
chosen generators of the same bidegree.
The quadratic and cubic generators in the following pictures have
homological degrees one and two, respectively.

\begin{example}[The strand $F_{18,2}$]\label{ex:original18}
The $\kk$-complex $F_{18,2}$ is supported on the labeled interval below:
\begin{center}
\begin{tikzpicture}[scale=0.8]
\node[circle,fill=black,inner sep=1.5pt,label=below:$x_5x_{13}$]
  (e1) at (-2,0) {};
\node[circle,fill=black,inner sep=1.5pt,label=below:$x_7x_{11}$]
  (e3) at (2,0) {};
\draw[thick,mid arrow={Stealth[reversed,length=1.45mm,width=1mm]}] (e3)--(e1);
\node at (0,0.3) {$e_{18,2}$};
\end{tikzpicture}
\end{center}
In particular, in homological degree zero we have
\[
 F_{0,18,2}=R_{18,2}
 =\kk x_5x_{13}\oplus\kk x_7x_{11},
\]
where the two monomials correspond to the vertices. In homological
degree one, $F_{1,18,2}=\kk e_{18,2}$, and the differential is the
usual boundary of an oriented edge:
\[
 \partial(e_{18,2})=x_7x_{11}-x_5x_{13}.
\]
\end{example}

\Needspace{6cm}
\begin{example}[The strand $F_{42,2}$]\label{ex:original42}
The maximum-product anchor is $x_{19}x_{23}$, and $F_{42,2}$ is
supported on the following star.
% Original coordinates and edge geometry; the center and right leaf
% have exchanged labels to put the required anchor at the center.
\begin{center}
\begin{tikzpicture}[scale=0.8]
\node[circle,fill=black,inner sep=1.5pt,label=below:$x_{19}x_{23}$]
  (e1) at (0,0) {};
\node[circle,fill=black,inner sep=1.5pt,label=above:$x_{11}x_{31}$]
  (e2) at (-4,4) {};
\node[circle,fill=black,inner sep=1.5pt,label=above:$x_{13}x_{29}$]
  (e3) at (0,4) {};
\node[circle,fill=black,inner sep=1.5pt,label=above:$x_5x_{37}$]
  (e4) at (4,4) {};
\draw[thick,mid arrow={Stealth[reversed,length=1.45mm,width=1mm]}] (e1)--(e2);
\draw[thick,mid arrow={Stealth[reversed,length=1.45mm,width=1mm]}] (e1)--(e3);
\draw[thick,mid arrow={Stealth[reversed,length=1.45mm,width=1mm]}] (e1)--(e4);
\node at (-2.8,2) {$e_{42,2}'$};
\node at (-0.5,2.7) {$e_{42,2}$};
\node at (2.8,2) {$e_{42,2}''$};
\end{tikzpicture}
\end{center}
The boundaries are
\[
\begin{aligned}
 \partial(e_{42,2})&=x_{19}x_{23}-x_{13}x_{29},\\
 \partial(e_{42,2}')&=x_{19}x_{23}-x_{11}x_{31},\\
 \partial(e_{42,2}'')&=x_{19}x_{23}-x_5x_{37}.
\end{aligned}
\]
In $\Gamma_{42,2}$ a single edge joins $x_{13}x_{29}$ to $x_{11}x_{31}$
because $13+29=11+31$. In the pictured star each edge is one basis
element. The path between these leaves goes first along $e_{42,2}$ to
the anchor, then backwards along $e_{42,2}'$. Its boundary is
\[
 \partial(e_{42,2}-e_{42,2}')=x_{11}x_{31}-x_{13}x_{29}.
\]
\end{example}

\Needspace{8cm}
\begin{example}[The strand $F_{23,3}$]\label{ex:original23}
The strand $F_{23,3}$ is supported on the labeled simplicial complex:
\nopagebreak
\begin{center}
\begin{tikzpicture}[scale=0.65,every node/.style={font=\small}]
\draw[fill=blue!12] (3,0)--(6,4)--(6,-4)--cycle;
\node[circle,fill=black,inner sep=1.5pt,label=below:$x_3^2x_{17}$]
  (e2) at (-3,0) {};
\node[circle,fill=black,inner sep=1.5pt,label=below left:$x_3x_7x_{13}$]
  (e4) at (3,0) {};
\node[circle,fill=black,inner sep=1.5pt,label=above:$x_5^2x_{13}$]
  (e6) at (6,4) {};
\node[circle,fill=black,inner sep=1.5pt,label=below:$x_5x_7x_{11}$]
  (e8) at (6,-4) {};
\draw[thick,mid arrow={Stealth[reversed,length=1.45mm,width=1mm]}] (e4)--(e2);
\draw[thick,mid arrow={Stealth[reversed,length=1.45mm,width=1mm]}] (e6)--(e4);
\draw[thick,mid arrow={Stealth[reversed,length=1.45mm,width=1mm]}] (e8)--(e6);
\draw[thick,mid arrow={Stealth[reversed,length=1.45mm,width=1mm]}] (e8)--(e4);
\node at (0,0.4) {$x_3e_{20,2}$};
\node at (3.7,2.3) {$x_{13}e_{10,2}$};
\node[right] at (6.2,0) {$x_5e_{18,2}$};
\node at (3.8,-2.2) {$x_7e_{16,2}$};
\node at (4.7,0) {$e_{23,3}$};
\end{tikzpicture}
\end{center}
The quadratic boundaries are
\[
\begin{aligned}
 \partial e_{10,2}&=x_5^2-x_3x_7,&
 \partial e_{16,2}&=x_5x_{11}-x_3x_{13},\\
 \partial e_{18,2}&=x_7x_{11}-x_5x_{13},&
 \partial e_{20,2}&=x_7x_{13}-x_3x_{17}.
\end{aligned}
\]
For instance, in $G_{23,3}$ we have
\[
 x_3x_7x_{13}=x_5^2x_{13}=x_5x_7x_{11}=x_3x_7x_{13}.
\]
This computation is visualized by tracing the boundary of the
two-dimensional cell labeled $e_{23,3}$. The differential is a cellular
boundary map with the additional requirement of $R$-linearity; for
example,
\[
 \partial(x_{13}e_{10,2})=x_{13}\partial(e_{10,2})
 =x_{13}(x_5^2-x_3x_7).
\]
In homological degree two, choose
\[
 \partial(e_{23,3})=x_5e_{18,2}-x_7e_{16,2}+x_{13}e_{10,2}.
\]
Applying $\partial$ again makes the three monomial differences
cancel. The graph has four vertices and four edges and is connected,
so this cycle spans its one-dimensional cycle space. Consequently
it supplies the entire second homological term in this strand over
every field.

The picture also illustrates variable truncation. All its labels
belong to $R^{<18}$. Passing to $R^{<14}$ removes the vertex
$x_3^2x_{17}$ and its incident edge $x_3e_{20,2}$, leaving the filled
triangle. Thus the cubic generator $e_{23,3}$ is already visible before
the full strand stabilizes. With these compatible choices of bases,
$F_{23,3}^{<m}=F_{23,3}$ for $m\geq17$. The relation cutoff
$F_{23,3}^{\le14}$ has the same three vertices but only the edge
$x_{13}e_{10,2}$. At $m=16$ it also has $x_7e_{16,2}$, giving a
two-edge path. The vertex $x_3^2x_{17}$ enters at $m=17$. At $m=18$ the
triangle fills in, while that vertex remains isolated until
$x_3e_{20,2}$ joins it to the triangle at $m=20$.
\end{example}

\Needspace{11cm}
\begin{example}[The strand $F_{31,3}$]\label{ex:original31}
There are six monomials in $R_{31,3}$. To specify the quadratic
generators used in this picture, set
\[
\begin{aligned}
 \partial e_{14,2}&=x_7^2-x_3x_{11},&
 \partial e_{24,2}&=x_{11}x_{13}-x_7x_{17},\\
 \partial e_{24,2}'&=x_{11}x_{13}-x_5x_{19},&
 \partial e_{26,2}&=x_{13}^2-x_7x_{19},\\
 \partial e_{26,2}'&=x_{13}^2-x_3x_{23},&
 \partial e_{28,2}&=x_{11}x_{17}-x_5x_{23}.
\end{aligned}
\]
Together with $e_{18,2}$ and $e_{20,2}$ above, these give the eight
edges below. Each edge carries its full generator label $x_pe_{s,2}^{(j)}$,
where $x_p$ is the fixed variable of the binary move. The middle quadrilateral and the two triangles
are the three two-dimensional cells.
\begin{center}
% The three faces share the two upper sides of the central quadrilateral.
\begin{tikzpicture}[scale=.96,line join=round,line cap=round,
    dot/.style={circle,fill=black,inner sep=1.5pt},
    vertex/.style={font=\normalsize,inner sep=2pt},
    edge label/.style={font=\small,inner sep=3pt},
    face label/.style={font=\normalsize,inner sep=2pt},
    ed/.style={semithick,mid arrow={Stealth[length=1.45mm,width=1mm]}}]
  \coordinate (F) at (0,7.0);
  \coordinate (D) at (-5.05,3.15);
  \coordinate (E) at (-1.8,3.65);
  \coordinate (B) at ( 1.8,3.65);
  \coordinate (C) at ( 5.05,4.35);
  \coordinate (A) at (0,.35);
  \fill[blue!12] (F)--(D)--(E)--(A)--(B)--(C)--cycle;
  % Orient each edge from the smaller product to the larger one.
  \draw[ed] (D)--(F);
  \draw[ed] (D)--(E);
  \draw[ed] (E)--(F);
  \draw[ed] (A)--(E);
  \draw[ed] (A)--(B);
  \draw[ed] (B)--(F);
  \draw[ed] (B)--(C);
  \draw[ed] (C)--(F);
  \path (D)--(F) node[midway,sloped,above,edge label] {$x_7e'_{24,2}$};
  \path (D)--(E) node[midway,sloped,below,edge label] {$x_5e_{26,2}$};
  \path (E)--(F) node[pos=.60,sloped,below,edge label] {$x_{13}e_{18,2}$};
  \path (A)--(E) node[midway,sloped,below,edge label] {$x_5e'_{26,2}$};
  \path (A)--(B) node[midway,sloped,below,edge label] {$x_3e_{28,2}$};
  \path (F)--(B) node[pos=.40,sloped,below,edge label] {$x_{11}e_{20,2}$};
  \path (B)--(C) node[midway,sloped,below,edge label] {$x_{17}e_{14,2}$};
  \path (F)--(C) node[midway,sloped,above,edge label] {$x_7e_{24,2}$};
  \foreach \name in {A,B,C,D,E,F} \node[dot] at (\name) {};
  \node[vertex,above=3pt] at (F) {$x_7x_{11}x_{13}$};
  \node[vertex,left=3pt] at (D) {$x_5x_7x_{19}$};
  \node[vertex,right=4pt] at (E) {$x_5x_{13}^2$};
  \node[vertex,left=4pt] at (B) {$x_3x_{11}x_{17}$};
  \node[vertex,right=3pt] at (C) {$x_7^2x_{17}$};
  \node[vertex,below=3pt] at (A) {$x_3x_5x_{23}$};
  \node[face label] at (-2.35,4.62) {$e'_{31,3}$};
  \node[face label] at (0,2.55) {$e_{31,3}$};
  \node[face label] at (2.35,4.99) {$e''_{31,3}$};
\end{tikzpicture}
\end{center}
The height in this drawing records the direction of the chosen
moves. Its top vertex is $x_7x_{11}x_{13}$, and its two sources are
$x_5x_7x_{19}$ and $x_3x_5x_{23}$. These are minimal elements
of the reachability order for the selected star edges. The numerical
prime-product order is total and has only the latter as its minimum.
The reachability order is not a lattice: the two sources have no common
lower bound. Nor is the displayed graph a Hasse diagram; for example,
its direct edge from $x_5x_7x_{19}$ to $x_7x_{11}x_{13}$ is
accompanied by a two-edge path through $x_5x_{13}^2$.

On the edge from $x_5x_7x_{19}$ to $x_7x_{11}x_{13}$, the fixed
factor $x_7$ moves from the second position to the first in the
sorted list of primes. The same change occurs for $x_5$ on the edge
from $x_3x_5x_{23}$ to $x_5x_{13}^2$. By contrast, on the edge
from $x_5x_{13}^2$ to $x_7x_{11}x_{13}$, the fixed $x_{13}$ can
occupy the third position at both ends.

The signed cellular boundaries are
\[
\begin{aligned}
 \partial e_{31,3}
   &=x_3e_{28,2}+x_{11}e_{20,2}-x_{13}e_{18,2}-x_5e_{26,2}',\\
 \partial e_{31,3}'
   &=x_5e_{26,2}-x_7e_{24,2}'+x_{13}e_{18,2},\\
 \partial e_{31,3}''
   &=x_{17}e_{14,2}+x_7e_{24,2}-x_{11}e_{20,2}.
\end{aligned}
\]
Each expression has boundary zero. They are independent because the
edges $x_3e_{28,2}$, $x_5e_{26,2}$, and $x_{17}e_{14,2}$ occur in just
one of the three respective expressions. The graph is connected,
with six vertices and eight edges, so its cycle space has dimension
three. Hence these are a basis of the cubic syzygies in this
bidegree, and the displayed complex supports $F_{31,3}$ over every
field.

The two computations
\[
\begin{aligned}
 x_3x_5x_{23}&=x_3x_{11}x_{17}=x_7^2x_{17},\\
 x_3x_5x_{23}&=x_5x_{13}^2=x_7x_{11}x_{13}=x_7^2x_{17}
\end{aligned}
\]
trace different paths between the same vertices. Their difference
is the boundary of the middle quadrilateral together with the right-hand triangle,
with compatible orientations. This is how a relation among sequences
of Goldbach moves appears in the next homological degree.
\end{example}


The last picture also realizes the compatible cutoff construction.
Its right-hand triangle, with generator $e_{31,3}''$, is present in
$F_{31,3}^{<17}$. At cutoff $19$ the left-hand triangle, with generator
$e_{31,3}'$, is added. At cutoff $23$ the middle quadrilateral, with generator
$e_{31,3}$, is added. Each stage retains every earlier cell and
attaching map. Thus these particular illustrated cutoffs form actual
polyhedral subcomplexes. This observation concerns the stated examples;
it does not assert such models for all strands or all cutoffs.

The next cubic example, in weight $33$, also connects these
polygonal strands to explicit Koszul classes. We give its picture
and calculation in Example~\ref{ex:weight33}, after defining the
Koszul complex.

\subsubsection{The quartic three-cell}
\begin{example}[The three-cell in $F_{34,4}$]
\label{ex:pyramid34}
The quadratic prime-product anchors and the following compatible
cubic choices give a unique three-cell in this strand. Its closure,
labeled $e_{34,4}$, is a square pyramid.
The diagram below is its top-down boundary view: the central
vertex is the apex, the four interior triangles are its lateral faces,
and the outer face represents its square base. Vertices carry monomial
labels, edges carry their quadratic-generator labels, and faces carry
their cubic-generator labels. The small edge arrows increase prime product.

\begin{center}
% Top-down boundary diagram of the square-pyramid three-cell in F_{34,4}.
% Labels and orientations checked by the illustrated-strands supplement.
% The outer face represents the square base; no interior three-cell label.
\begin{tikzpicture}[
  scale=1.05,line join=round,
  v/.style={circle,fill=black,inner sep=1.6pt},
  vertex label/.style={font=\normalsize,inner sep=3pt},
  edge label/.style={font=\small,inner sep=2pt},
  face label/.style={font=\normalsize,inner sep=2pt},
  ed/.style={line width=.7pt,mid arrow={Stealth[reversed,length=1.45mm,width=1mm]}}]
\coordinate (A) at (-4,-4);
\coordinate (B) at (4,-4);
\coordinate (C) at (4,4);
\coordinate (D) at (-4,4);
\coordinate (U) at (0,0);
\fill[blue!12] (A)--(B)--(C)--(D)--cycle;
% The four outer edges.
\draw[ed] (B)--(A);
\draw[ed] (C)--(B);
\draw[ed] (C)--(D);
\draw[ed] (A)--(D);
\node[edge label,below=7pt] at (0,-4) {$x_5x_{13}e_{16,2}$};
\node[edge label,right=7pt] at (4,0) {$x_5x_{11}e_{18,2}$};
\node[edge label,above=7pt] at (0,4) {$x_5x_7e_{22,2}'$};
\node[edge label,left=7pt] at (-4,0) {$x_3x_5e_{26,2}$};
% The four edges incident with the apex, with full quadratic-generator labels.
\draw[ed] (B)--(U)
  node[pos=.53,sloped,below=4pt,edge label] {$x_{11}x_{13}e_{10,2}$};
\draw[ed] (U)--(A)
  node[pos=.53,sloped,above=4pt,edge label] {$x_3x_{13}e_{18,2}$};
\draw[ed] (C)--(U)
  node[pos=.53,sloped,above=4pt,edge label] {$x_7x_{11}e_{16,2}$};
\draw[ed] (U)--(D)
  node[pos=.53,sloped,above=4pt,edge label] {$x_3x_7e'_{24,2}$};
% Monomial vertices. A small leader identifies the central monomial cleanly.
\foreach \P in {A,B,C,D,U}{\node[v] at (\P) {};}
\node[vertex label,below=5pt] at (A) {$x_3x_5x_{13}^{2}$};
\node[vertex label,below=5pt] at (B) {$x_5^{2}x_{11}x_{13}$};
\node[vertex label,above=5pt] at (C) {$x_5x_7x_{11}^{2}$};
\node[vertex label,above=5pt] at (D) {$x_3x_5x_7x_{19}$};
\node[vertex label] (apex label) at (0,1.25) {$x_3x_7x_{11}x_{13}$};
\draw[-{Stealth[length=1.2mm,width=.8mm]},line width=.4pt]
  (apex label.south)--(0,.12);
% Four triangular two-cells and the outer square two-cell.
\node[face label] at (0,-2.85) {$x_{13}e_{21,3}$};
\node[face label] at (2.85,0) {$x_{11}e_{23,3}$};
\node[face label] at (0,2.85) {$x_7e_{27,3}$};
\node[face label] at (-2.85,0) {$x_3e'_{31,3}$};
\node[face label] at (0,5.05) {Outer face: $x_5e'_{29,3}$};
\end{tikzpicture}

\end{center}
To specify the remaining boundaries, set
$\partial e_{22,2}'=x_{11}^2-x_3x_{19}$ and choose
\[
\begin{aligned}
 \partial e_{21,3}
   &=x_{11}e_{10,2}-x_5e_{16,2}+x_3e_{18,2},\\
 \partial e_{27,3}
   &=x_{11}e_{16,2}-x_5e_{22,2}'+x_3e'_{24,2},\\
 \partial e'_{29,3}
   &=x_{13}e_{16,2}+x_{11}e_{18,2}
      -x_7e_{22,2}'+x_3e_{26,2}.
\end{aligned}
\]
The boundaries of $e_{23,3}$ and $e'_{31,3}$ were given above. With compatible
face orientations, the new generator has boundary
\[
 \partial e_{34,4}
 =x_{13}e_{21,3}-x_{11}e_{23,3}-x_7e_{27,3}
   +x_5e'_{29,3}-x_3e'_{31,3}.
\]
The left triangular face is $x_3e'_{31,3}$, obtained by multiplying
the left triangle of Example~\ref{ex:original31} by $x_3$.
The square base is labeled $x_5e'_{29,3}$.

Every edge occurs twice with opposite incidence signs. The resulting
two-sphere uses five vertices, eight edges, and five faces; filling it
gives the three-cell. The full strand has dimensions $(11,20,11,1)$,
so this picture shows the support of its unique three-cell rather
than the entire strand.

For completeness, the finite integral checks for these examples are
summarized below. The dimensions refer to the full resolution strand,
including monomial multiples of lower generators.
\[
\begin{array}{c|c|c}
 (n,d)&(\dim F_{0,n,d},\dim F_{1,n,d},\ldots)
       &(\operatorname{rank}\partial_1,\operatorname{rank}\partial_2,\ldots)\\\hline
 (18,2)&(2,1)&(1)\\
 (42,2)&(4,3)&(3)\\
 (23,3)&(4,4,1)&(3,1)\\
 (31,3)&(6,8,3)&(5,3)\\
 (33,3)&(8,11,4)&(7,4)\\
 (34,4)&(11,20,11,1)&(10,10,1)
\end{array}
\]
Here is an exact way to verify the last line, rather than infer
uniqueness from the pictured sphere. Enumerate the eleven quartic
monomials, all twenty monomial multiples of the quadratic star
relations, and all eleven variable multiples of the chosen cubic
cycles. Their integer boundary matrices have consecutive products
zero. Integer elimination uses only pivots $1$ or $-1$ and leaves
zero residual matrices, with ranks $10$ and $10$ for the first two
maps. The displayed three-cell boundary is a primitive generator
of the remaining one-dimensional kernel. Thus the existing cubic
multiples fill the entire first cycle group, and exactly one new
third generator is required. These checks hold over every field.

The data are in
\path{supplement/computations/illustrated-strands-data.json}, with the
standard-library checker
\path{supplement/computations/illustrated-strands-check.py}. It reconstructs the
matrices from polynomial boundaries, checks all cubic generator
weights through $31$, and verifies the five face boundaries of this
pyramid. For the example in weight $33$, \path{supplement/computations/koszul33-check.py}
checks the four polygon cycles, the swaps, and their Koszul classes over the integers. All polynomial
coefficients in positive differentials lie in $\m$, so these are
minimal generators. The finite choices can be extended to a minimal
resolution by the preceding construction. This verifies the examples
without assuming a global cellular resolution.
\end{example}

The companion paper \cite{CompanionTopology} compares these binary
strands with strands for the ternary relation ideal.

The binary examples above show how minimal binary moves form the first stage of
an inductive construction: polygon boundaries supply cubic syzygies,
and compatible polygon faces can bound a three-cell. In general,
however, the next step chooses minimal homogeneous generators of a
syzygy module. A product order on monomials does not yet choose those
higher generators or ensure that their boundaries are spheres.
Extending the order and the moves to compatible higher cells is an
open problem here; the displayed models establish only the stated
strands and cutoff inclusions.
